% Use from this directory, or adjust image paths when copying into the paper.
% Requires graphicx and amsmath. Figures target 7.6-inch full text width.

\begin{figure*}[t]
\centering
\includegraphics[width=\textwidth]{Figure_B_neighborhood_merged.pdf}
\caption{\textbf{Task0 neighborhood diagnostics for EWCLoRA and FS-EWCLoRA.} Existing ImageNet-A seed0 checkpoints are evaluated with deterministic transforms and each method's fixed 20-class Task0 head. Radius $r$ is the joint Frobenius distance from its own Task0 anchor across 24 effective Q/V weight matrices. The same $K=16$ unit directions are shared across methods and positive radii; $r=0$ reuses one origin per method. For each direction, positive damage is $[L(w_0+ru)-L(w_0)]_+$, using mean cross-entropy on 1,259 Task0 training or 307 Task0 test images. Panels (a) and (b) show its directional mean $G_{\mathrm{mean}}$ and empirical 95th percentile $G_{95}$, respectively. Panel (c) shows the mean and $Q_{95}$ of raw empirical Fisher traces across directions, estimated on a fixed subset of 256 training images without normalization, flooring, or regularization. Quantiles use linear interpolation. Panel (d) gives the prefix-safe tolerance radius $\hat r_\epsilon$, the largest sampled radius for which $G_{\mathrm{mean}}\leq\epsilon$ at every sampled radius up to that point. Triangles mark right-censoring at the maximum probed radius, 20; nested markers preserve exact method ties. Solid/filled and dashed/open denote train and test in (a), (b), and (d). Radius axes in (a--c) and damage axes in (a--b) use symmetric-log scales with linear thresholds 0.5 and $10^{-4}$; (d) uses logarithmic $\epsilon$. Connecting lines guide the eye; these finite-direction summaries are not confidence intervals or worst-case guarantees.}
\label{fig:B-neighborhood-merged}
\end{figure*}

\begin{figure*}[t]
\centering
\includegraphics[width=\textwidth]{Figure_C_sequence_row.pdf}
\caption{Existing ImageNet-A sequences (seed 0), evaluated in full FP32. Learning stage $T=t_{\mathrm{code}}+1$ counts learned tasks: $T=1$ is the Task0 anchor and $T=10$ follows Task9. EWCLoRA uses SGD at Task0 and SGD+EWC thereafter; FS-EWCLoRA uses GAM at Task0 and GAM+EWC thereafter, so subsequent optimizers differ. (a) Actual displacement $r_T=\|W_{T-1}-W_0\|_2$ from each method's own anchor, in the joint coordinates of 24 effective Q/V $BA$ matrices (scaling 1). (b) Signed Task0 test cross-entropy change $D_T=L_0^{\mathrm{test}}(W_{T-1})-L_0^{\mathrm{test}}(W_0)$; negative values indicate improvement. (c) Task0 test accuracy. Both test metrics use the same 307 examples and fixed 20-class Task0 head. Lines join all measured checkpoints; negative and nonmonotonic changes are retained. (d) Final chord-normalized sensitivity $\bar S_{10}=I_{10}/r_{10}^2$, with $I_{10}=\int_0^1\Delta W_{10}^{\top}\widehat F_{\mathrm{emp}}(W_0+s\Delta W_{10})\Delta W_{10}\,ds$ and $\Delta W_{10}=W_9-W_0$. The raw empirical Fisher averages outer products of unscaled per-example CE gradients on 256 stratified Task0 training examples, using the fixed head without trace normalization, flooring, or regularization. Nine composite-Simpson nodes evaluate only the final straight effective-weight chord, not the optimizer trajectory. This single-run comparison supplies no across-run uncertainty or isolated effect of the Task0 anchor; displacement is not a damage-tolerance radius.}
\label{fig:sequence}
\end{figure*}
